\chapter{Conclusion}
\label{sec:conclusion}

\glsresetall
This thesis compared three approaches to \gls{ats} at sentence and document granularity: an off-the-shelf simplifier (the \emph{comparison simplifier}), a \gls{seqseq} model fine-tuned for this thesis, and prompted open-weight \glspl{llm}.
All three are served by a browser extension that simplifies web pages in place, and this deployment was audited on twelve authentic pages.

\textbf{\Cref{rq:comparison}} asks how the approaches compare.
No approach wins outright: the ranking depends on which systems are compared and on which metric is used.
Fine-tuning the 139\gls{million}-parameter base version of \gls{bart} on a task-matched corpus clearly improves on the untuned checkpoint on \gls{sari} and \gls{dsari}.
\Gls{sari} rises from \num{21.39} to \num{37.80} at sentence level (\(+16.42\), \glsxtrlong{ci} \([+15.38, +17.50]\), \(\gls{pval} < 0.001\)), and \gls{dsari} from \num{14.34} to \num{35.44} on the complete document test split (\(\gls{pval} < 0.001\)).
These are the only comparisons tested for significance.
Against the comparison simplifier, which is three times its size, the fine-tuned checkpoint scores \num{0.37}~\gls{sari} lower, well inside the \({\approx}2\)-point noise floor, so its advantage over that model is serving cost.

Both prompted \glspl{llm} score higher than the fine-tuned checkpoint on \gls{sari} at sentence level.
At document level, the prompted 7\gls{billion}-parameter model scores higher on \gls{lens} but lower on \gls{dsari}, consistently across three sampling seeds.
\Gls{lens} cannot settle this disagreement, because it is applied outside the sentence-level data it was trained on \parencite[16385--16386]{maddela-etal-2023-lens}, so the document-level ordering remains open.

Structured manual inspection found errors that the aggregate scores did not reveal: a document checkpoint trained on less data changed the census year 2010 to 2000 in 27 of 35 affected documents.

\textbf{\Cref{rq:deployment}} asks which content-selection and replacement strategies make a deployed simplifier reliable on real pages, and what an audit shows about its robustness.
The exploratory site audit gives a mixed answer: replacement is reliable, content selection mostly is, and most emptied links trace back to the model's edits rather than to the replacement.
Replacement preserved the page's \gls{dom} on all 24 runs (twelve pages at both granularities): no tag, image source or link target was lost, and reverting restored the page's text and attributes every time.
Content selection mostly withholds the right text, but its heuristic for code-like text also rejects ordinary prose that contains a semicolon.

The most serious result of the audit is that five sentences were served with a negation deleted, which reversed their meaning, and no guard fired.
Prompt injection also succeeded end to end on both prompted \glspl{llm}.
Whether the prompts' audience parameter serves the audiences it names was not tested.

This thesis makes four contributions (\cref{sec:contributions}).
It assesses the corpora used and documents data-quality issues in them, such as encoding errors and overlap between training and test data.
It compares the three approaches at both granularities in one evaluation harness, including the complete \num{8000}-document D-Wikipedia test split.
It provides a \gls{dom}-preserving browser extension in which the approaches can be compared on authentic web pages.
Finally, it derives methodological findings about the evaluation of \gls{ats} (\cref{sec:cross-cutting}).
These are three observations that plausibly transfer beyond the models used here: the unit a model is trained on should match the unit a deployment feeds it; the choice of benchmark and metric can change a conclusion; and inspecting outputs finds failures that the scores do not show.

I consider the deployment half of this thesis the more useful one.
Benchmark scores said nothing about the negation deletions, the withheld prose or the sub-sentence fragments sent to the model; each was found only on closer inspection of the deployed system.
In my view, work of that kind is under-reported in \gls{ats} research relative to what it costs to find.
Since 2010, that research has been dominated by sentence-aligned Simple English Wikipedia data \parencite[402]{xu-etal-2016-optimizing}.

The next steps follow from what remains open (\cref{sec:future-work}): a human study to settle the document-level ordering, a polarity check so that the change guard catches deleted negations, and a language gate before any multilingual training.

The motivating example of \cref{sec:motivation} shows why benchmark scores alone are not enough.
The sentence-level \gls{seqseq} models in this thesis turned the four-word heading ``Municipal recycling policy revisions'' into ``Other websites''.
The output may originate in the training corpus, which maps short fragments with little content onto exactly that phrase in 38 pairs.
A system pointed at a live page receives inputs whose handling no benchmark score describes.

Finding such failures took something other than a better metric.
It took a log of what the model produced and what was served; a harness that runs the extension's own code rather than a reimplementation of it; and structured inspection of outputs.
The failures and findings that mattered most in this deployment did not show up in the benchmark scores.
